\documentclass[12pt,a4paper]{article}
\usepackage[UTF8]{ctex}
\usepackage{geometry}
\usepackage{amsmath,amssymb,amsthm}
\usepackage{graphicx}
\usepackage{booktabs}
\usepackage{longtable}
\usepackage{array}
\usepackage{enumitem}
\usepackage{hyperref}
\usepackage{fancyhdr}
\usepackage{xcolor}
\usepackage{listings}
\usepackage{setspace}  % 添加行距控制包

\geometry{left=2.5cm,right=2.5cm,top=2.5cm,bottom=2.5cm}
\setstretch{1.5}  % 设置1.5倍行距（CNIPA推荐）

\pagestyle{fancy}
\fancyhf{}
\fancyhead[C]{发明专利技术交底书}
\fancyfoot[C]{\thepage}

\title{\textbf{发明专利技术交底书} \\[0.5em]
\Large 面向 5G/6G 核心网与跨链桥的基于 Falcon 算法的量子安全门限签名系统及方法}
\author{中国移动通信有限公司研究院}
\date{\today}

\begin{document}

\maketitle

\tableofcontents
\newpage

%==============================================================================
\section{发明名称}
%==============================================================================

面向 5G/6G 核心网与区块链跨链桥的基于 Falcon 算法的量子安全门限签名系统及方法

\textbf{英文名称}: Quantum-Secure Threshold Signature System and Method Based on Falcon Algorithm for 5G/6G Core Network and Blockchain Cross-Chain Bridge Scenarios

%==============================================================================
\section{申请人信息}
%==============================================================================

\begin{table}[h]
\centering
\begin{tabular}{|p{4cm}|p{8cm}|}
\hline
\textbf{字段} & \textbf{内容} \\
\hline
申请人名称 & 中国移动通信有限公司研究院 \\
\hline
申请人类型 & 企业 \\
\hline
国籍/注册地 & 中国 \\
\hline
统一社会信用代码 & 9111010279160212X8 \\
\hline
通讯地址 & 北京市西城区宣武门西大街32号 移动创新大楼 \\
\hline
邮政编码 & 100053 \\
\hline
联系电话 & +86-13521894156 \\
\hline
电子邮箱 & xudayj@chinamobile.com \\
\hline
\end{tabular}
\end{table}

%==============================================================================
\section{发明人信息}
%==============================================================================

\subsection{第一发明人}

\begin{table}[h]
\centering
\begin{tabular}{|p{4cm}|p{8cm}|}
\hline
\textbf{字段} & \textbf{内容} \\
\hline
姓名 & 许达 (Xu Da) \\
\hline
国籍 & 美国 \\
\hline
工作单位 & 中国移动通信有限公司研究院 \\
\hline
职务/职称 & 主任研究员 \\
\hline
通讯地址 & 北京市西城区宣武门西大街32号 移动创新大楼 \\
\hline
邮政编码 & 100053 \\
\hline
联系电话 & +86-13521894156 \\
\hline
电子邮箱 & xudayj@chinamobile.com \\
\hline
\end{tabular}
\end{table}

%==============================================================================
\section{技术领域}
%==============================================================================

本发明涉及移动通信网络安全、区块链跨链互操作、后量子密码学（Post-Quantum Cryptography, PQC）及多方安全计算（Multi-Party Computation, MPC）领域，特别涉及一种部署于 5G/6G 移动核心网网元（AMF/UDM/AUSF/NRF）协同认证场景以及区块链跨链桥多方托管场景下，由分布式节点协同生成符合 NIST FIPS 206 标准 Falcon 算法数字签名的技术系统及方法。本发明并非一种独立的纯算法，而是与上述两类具体通信与金融基础设施场景的网元部署形态、信令交互流程、链上验证合约及容灾切换机制深度耦合的工程化技术方案。

%==============================================================================
\section{背景技术}
%==============================================================================

\subsection{现有技术的局限性}

本发明针对的两类具体场景在现网部署中均面临相同的量子安全短板，但其威胁路径与改造约束截然不同：

\textbf{场景一：5G/6G 核心网网元间的协同认证与切片签名}。3GPP TS 33.501/TS 33.310 规定 AMF、AUSF、UDM、NRF、SEPP 等核心网网元在 SBA（服务化架构）总线上交互时，必须使用基于 X.509 与 TLS 1.3 的相互认证；其证书签发与切片专用网元间的事件订阅签名当前均默认采用 ECDSA P-256/P-384 或 EdDSA。在 N32-c/N32-f 漫游接口、网络切片实例化（NSI）创建消息以及面向行业专网的能力开放接口（NEF）上，运营商每天产生 $10^7$ 量级的签名动作，一旦其中长期密钥被量子计算机攻破，将导致跨运营商漫游伪造、切片越权与计费篡改。

\textbf{场景二：跨链桥多方托管下的资产签名}。主流跨链桥（Wormhole、LayerZero、Axelar 等）由 $n$ 个守护者节点对跨链消息进行 $(t,n)$ 门限签名，承载资产规模数十亿美元，普遍采用 ECDSA、EdDSA 或 BLS 聚合签名：

\begin{enumerate}
    \item \textbf{ECDSA 门限签名}：基于椭圆曲线离散对数问题
    \item \textbf{EdDSA 门限签名}：基于 Edwards 曲线的签名方案
    \item \textbf{BLS 聚合签名}：基于双线性配对的签名方案
\end{enumerate}

上述两类场景中所使用的椭圆曲线签名方案均面临量子计算机的威胁；并且简单地将单机版 NIST PQC 算法（Dilithium/Falcon/Sphincs+）直接部署到上述场景，会引入新的工程难题：单点持有完整私钥不符合电信运营商内控合规与跨链桥多托管原则，必须以门限/MPC 形式部署；而 Falcon 等格基方案的浮点高斯采样、拒绝采样、NTT 多项式乘法等操作直接做 MPC 化会导致通信轮数爆炸，无法满足 5G 核心网毫秒级 SLA 与跨链桥分钟级清算窗口的实际约束。

\subsection{量子攻击形式化分析}

\textbf{定理（Shor 算法复杂度）}：给定 $n$ 位整数 $N$ 或阶约为 $2^n$ 的椭圆曲线群，Shor 算法分解 $N$ 或求解离散对数问题的时间复杂度为：
$$T_{\text{quantum}} = O(n^3) \text{ 次量子操作}$$
需要 $O(n)$ 个逻辑量子比特。

\textbf{推论（ECDSA 脆弱性）}：对于 256 位 ECDSA (secp256k1)，一台拥有约 2,330 个逻辑量子比特的量子计算机可以在 $O(256^3) \approx 2^{24}$ 次量子操作内破解该方案，使其安全性归零。

\textbf{Grover 算法考量}：虽然 Grover 算法对搜索问题仅提供二次加速：
$$T_{\text{Grover}} = O(\sqrt{2^n}) = O(2^{n/2})$$
但这影响签名中使用的对称原语。本设计使用 256 位输出的 SHA-3/SHAKE，提供：
\begin{align}
\text{后量子抗碰撞性} &= 256/2 = 128 \text{ 位（通过生日界）} \\
\text{后量子抗原像性} &= 256/2 = 128 \text{ 位（通过 Grover）}
\end{align}

\textbf{时间紧迫性}：NIST 估计具备密码学相关能力的量子计算机可能在 2030-2035 年出现。管理数十亿资产的跨链桥现在就需要进行量子安全迁移。

\subsection{Falcon 算法及其分布式实现的挑战}

Falcon（Fast-Fourier Lattice-based Compact Signatures over NTRU）是 NIST 于 2022 年选定的三种后量子数字签名标准之一（FIPS 204/205/206）。该算法具有以下优势：

\begin{itemize}
    \item \textbf{签名长度短}：约 666 字节 (Falcon-512) 至 1280 字节 (Falcon-1024)
    \item \textbf{验证速度快}：由于高效的 NTT 结构，比 Dilithium 快约 10 倍
    \item \textbf{基于格的安全性}：基于 NTRU/Ring-SIS 困难性，抵抗已知量子攻击
\end{itemize}

然而，Falcon 算法的分布式实现（即门限签名）面临重大技术挑战：

\textbf{挑战 1：MPC 中的离散高斯采样}

Falcon 签名的核心步骤需要在 NTRU 格上进行离散高斯分布采样：
$$\mathbf{s} \leftarrow D_{\mathbf{B}, \sigma, \mathbf{c}}$$
其中 $\mathbf{B}$ 是私钥陷门基，$\sigma$ 是高斯参数，$\mathbf{c}$ 是哈希导出的中心。

\textbf{挑战 2：基于 NTT 的分布式多项式运算}

Falcon 依赖数论变换（NTT）进行高效多项式乘法：
$$f \cdot g = \text{iNTT}(\text{NTT}(f) \odot \text{NTT}(g))$$

\textbf{挑战 3：拒绝采样通信开销}

Falcon 使用拒绝采样来确保签名的统计独立性（零泄露）：
$$\Pr[\text{accept}] = \frac{1}{M} \cdot \exp\left(-\frac{\langle \mathbf{s}, \mathbf{c} \rangle}{\sigma^2}\right) \approx 0.65$$

\textbf{挑战 4：昂贵的链上验证成本}

\begin{itemize}
    \item \textbf{Dilithium}：签名较大（约2.4KB）导致高存储和 Gas 成本（以太坊上约 180,000 Gas）
    \item \textbf{Sphincs+}：签名更大（约8KB），不适合频繁的跨链交易
    \item \textbf{Falcon}：提供最小的签名尺寸（约666 字节）和最低的验证成本（约50,000 Gas）
\end{itemize}

%==============================================================================
\section{发明内容}
%==============================================================================

\subsection{发明目的}

本发明旨在解决上述技术问题，提供一种与 5G/6G 核心网网元协同认证场景和区块链跨链桥多方托管场景深度耦合的量子安全门限 Falcon 签名系统及方法。本发明并非提供一种通用的抽象 MPC 算法，而是针对上述两类场景各自的网元形态、信令流程、SLA 约束与验证载体（SBA 总线 JOSE/COSE 报文 vs. EVM 兼容链上验证合约），提供可直接落地的系统架构、协议步骤、参数配置与容灾策略，并通过本机实测提供可复核的量化证据（见第\ref{sec:effect}节技术效果）。

\subsection{技术方案}

本发明提出以下核心创新点：

\subsubsection{创新点 A：基于 NTRU 结构的算术共享 NTT 协议}

设计了一种新型 MPC 协议，将 Falcon 实现中使用的线性变换表示分解为多个子计算。

设私钥多项式 $f$ 在 $N$ 个节点间进行算术共享为 $[f]_1, [f]_2, ..., [f]_N$，满足：
$$f = \sum_{i=1}^{N} [f]_i$$

对于选定的线性变换操作 $\mathcal{T}(\cdot)$，利用线性性质：
$$\mathcal{T}(f) = \mathcal{T}\left(\sum_{i=1}^{N} [f]_i\right) = \sum_{i=1}^{N} \mathcal{T}([f]_i)$$

每个节点独立计算其本地分片的变换，从而在不重构私钥的情况下实现全局变换域计算。

\subsubsection{创新点 B：具有正确参数校准的分布式高斯采样}

一项关键创新解决了正确分布的高斯样本的分布式生成问题。

\textbf{定理（高斯聚合）}：如果每一方 $P_i$ 采样 $[z]_i \leftarrow D_{\sigma/\sqrt{N}, R}$，则聚合值 $z = \sum_{i=1}^{N} [z]_i$ 服从分布 $D_{\sigma, R}$。

\textbf{证明}：对于独立高斯分布，方差相加：$\text{Var}(z) = \sum_{i=1}^{N} \text{Var}([z]_i) = N \cdot (\sigma/\sqrt{N})^2 = \sigma^2$。

\textbf{注}：严格来说，对于格上的离散高斯分布，只要缩放参数 $\sigma/\sqrt{N}$ 超过格的平滑参数 $\eta_\epsilon(\Lambda)$，上述卷积性质在统计距离上即成立。本发明的参数选择确保了满足此条件。

这要求每一方使用缩放参数 $\sigma_i = \sigma/\sqrt{N}$ 而非全局参数 $\sigma$。

\subsubsection{创新点 C：基于 Beaver 三元组预处理的安全分布式范数验证}

针对 Falcon 签名中的签名范数边界检查，发明了一种两阶段协议，确保在不泄露私钥分片的情况下验证全局签名范数是否满足 $\|\mathbf{s}\|^2 \le \lfloor \beta^2 \rfloor$：

\textbf{离线阶段（预处理）}：
\begin{itemize}
    \item 各方生成 Beaver 三元组 $([a]_i, [b]_i, [c]_i)$，其中 $c = a \cdot b$
    \item 这些三元组用于在线阶段高效计算交叉项 $\langle [s]_i, [s]_j \rangle$
\end{itemize}

\textbf{在线阶段（常数轮次）}：
\begin{enumerate}
    \item \textbf{本地范数计算}：每个节点 $P_i$ 计算本地范数 $\|[s]_i\|^2$ 并生成掩码 $m_i$
    \item \textbf{基于 Beaver 三元组的交叉项计算}：利用预处理的三元组，各方在 2 轮内计算：
    $$\sum_{i<j} \langle [s]_i, [s]_j \rangle$$
    \item \textbf{全局范数组装}：计算 $\|s\|^2 = \sum_i \|[s]_i\|^2 + 2\sum_{i<j} \langle [s]_i, [s]_j \rangle$
    \item \textbf{分布式边界检查}：各方验证是否 $\|s\|^2 \le \lfloor \beta^2 \rfloor$，若不满足则重启采样
    \item \textbf{条件揭示}：仅在验证通过时才揭示实际签名分量
\end{enumerate}

\textbf{通信复杂度}：在线 $O(1)$ 轮（每批签名需 $O(n^2)$ 离线预处理）。

\subsubsection{创新点 D：具有作弊检测的可验证秘密共享}

为确保针对恶意参与方的安全性，系统集成了：

\begin{enumerate}
    \item \textbf{用于密钥分片的 Feldman 风格 VSS}：在密钥生成期间，每一方 $P_i$ 发布承诺：
    $$C_i = g^{[f]_i} \mod p$$
    
    \item \textbf{基于承诺的作弊检测}：每个签名轮次以绑定承诺 $C_i = H(m_i \| [s]_i)$ 开始
    
    \item \textbf{中止并识别协议}：当验证失败时，识别并排除作弊方
    
    \item \textbf{针对移动对手的主动刷新}：定期密钥分片刷新
\end{enumerate}

\subsubsection{创新点 E：动态节点准入与密钥重构}

利用 NTRU 格的线性同态性质，实现了以下功能：

\begin{enumerate}
    \item \textbf{动态节点添加}：新节点加入时，现有节点通过秘密共享协议分配新的私钥分片，同时保持主公钥不变
    \item \textbf{节点撤销}：通过主动秘密共享更新剩余节点的私钥分片
    \item \textbf{自动修复}：检测到节点离线时，其他节点协同重构该节点的私钥分片
\end{enumerate}

\subsubsection{创新点 F：容错与超时处理}

协议包含针对实际部署的健壮错误处理机制：

\begin{itemize}
    \item \textbf{超时检测}：每个协议阶段具有可配置的超时时间
    \item \textbf{优雅降级}：参与方超时时，从当前签名尝试中排除
    \item \textbf{网络分区处理}：各方之间维持心跳
    \item \textbf{状态恢复}：每个参与方在每个阶段后持久化协议状态
\end{itemize}

%==============================================================================
\section{附图说明}
%==============================================================================

\textbf{图 1} 是展示源链、门限签名系统和目标链之间交互流程的整体系统架构示意图。

\begin{figure}[h]
    \centering
    \includegraphics[width=0.9\textwidth]{../Figure_1.png}
    \caption{系统架构示意图}
    \label{fig:arch}
\end{figure}

\textbf{图 2} 是说明安全分布式范数验证过程的流程图，包括承诺、预检查和揭示阶段。

\begin{figure}[h]
    \centering
    \includegraphics[width=0.9\textwidth]{../Figure_2.png}
    \caption{安全分布式范数验证流程图}
    \label{fig:sampling}
\end{figure}

\textbf{图 3} 是说明动态节点管理场景的示意图，包括节点加入、撤销和离线恢复。

\begin{figure}[h]
    \centering
    \includegraphics[width=0.9\textwidth]{../Figure_3.png}
    \caption{动态节点管理示意图}
    \label{fig:dynamic}
\end{figure}

%==============================================================================
\section{具体实施方式}
%==============================================================================

\subsection{系统架构}

本发明的系统架构针对两类落地场景采用同一底层签名引擎与差异化适配层。整体如下：

\begin{verbatim}
┌─────────────────────────────────────────────────────────────┐
│  应用层（场景适配）                                          │
│  ┌──────────────────────┐   ┌──────────────────────────┐    │
│  │ 5G/6G 核心网适配器   │   │  跨链桥适配器             │    │
│  │ - SBA/N32 报文封装   │   │ - EVM/非EVM 链消息封装    │    │
│  │ - JOSE/COSE 签名头   │   │ - Merkle 批处理           │    │
│  │ - 3GPP TS 33.310 兼容 │   │ - 守护者节点共识接入      │    │
│  └──────────────────────┘   └──────────────────────────┘    │
├─────────────────────────────────────────────────────────────┤
│  ┌──────────────┐  ┌──────────────┐  ┌──────────────┐       │
│  │   节点 P₁    │  │   节点 P₂    │  │   节点 Pₙ    │       │
│  │ 密钥分片[f]₁ │  │ 密钥分片[f]₂ │  │ 密钥分片[f]ₙ │       │
│  │   本地 NTT   │  │   本地 NTT   │  │   本地 NTT   │       │
│  │   掩码生成   │  │   掩码生成   │  │   掩码生成   │       │
│  └──────────────┘  └──────────────┘  └──────────────┘       │
├─────────────────────────────────────────────────────────────┤
│                MPC 协调层（通信）                            │
├─────────────────────────────────────────────────────────────┤
│  验证层：                                                    │
│   - 5G 场景：SEPP/NRF 内置 Falcon 验证模块                   │
│   - 跨链场景：EVM 智能合约 Falcon 验证（约 50,000 Gas）      │
└─────────────────────────────────────────────────────────────┘
\end{verbatim}

在 5G/6G 核心网场景中，签名节点 $P_i$ 部署为运营商核心网内的"PQC-Signing NF"网元，按 3GPP SBA 规范注册到 NRF；在跨链桥场景中，签名节点 $P_i$ 部署为跨链桥守护者节点（Guardian）的安全飞地进程，由跨链合约共识层触发签名请求。两种部署形态共享同一 D-KeyGen/D-Sign 内核，但适配层、SLA 配置、容灾策略与验证载体均针对场景定制，使本发明的技术方案与场景密不可分。

\subsection{实施步骤}

\subsubsection{实施例 1：核心协议流程}

\textbf{步骤 1：带 VSS 的分布式密钥生成 (D-KeyGen)}

\begin{enumerate}
    \item \textbf{初始化}：节点组 $(P_1, ..., P_n)$ 商定系统参数
    \item \textbf{分片生成}：每个节点 $P_i$ 从缩放的高斯分布 $D_{\sigma/\sqrt{N}, R}$ 中采样本地多项式分片 $[f]_i, [g]_i$
    \item \textbf{可验证秘密共享 (VSS)}：每个节点广播承诺 $C_i = \text{Commit}([f]_i)$
    \item \textbf{陷门计算}：节点运行 MPC-Extended-GCD 协议计算 $(F, G)$ 的分片
    \item \textbf{输出}：每个节点将 $([f]_i, [g]_i, [F]_i, [G]_i)$ 存储在安全存储器中
\end{enumerate}

\textbf{步骤 2：高性能离线预处理}

\begin{enumerate}
    \item \textbf{OT 扩展}：节点利用不经意传输扩展高效生成 OT
    \item \textbf{三元组生成}：利用 OT 生成 Beaver 三元组 $([a], [b], [c])$
    \item \textbf{正确性验证}：执行"牺牲"步骤验证三元组正确性
    \item \textbf{存储}：验证后的三元组存储在"三元组队列"中
\end{enumerate}

\textbf{步骤 3：消息映射与请求处理}

\begin{enumerate}
    \item \textbf{请求摄入}：系统接收跨链桥请求
    \item \textbf{去重}：共识层确保所有节点对请求序列达成一致
    \item \textbf{哈希}：节点本地计算 $c = H(r || M)$ 并映射到环 $R$
\end{enumerate}

\textbf{步骤 4：隐私保护分布式采样}

\begin{enumerate}
    \item \textbf{本地采样}：每个节点 $P_i$ 采样 $[z]_i \leftarrow D_{\sigma/\sqrt{N}, R}$
    \item \textbf{零知识证明（可选）}：附带非交互式零知识证明
    \item \textbf{分片计算}：$[s]_i = [t]_i + [z]_i$
\end{enumerate}

\textbf{步骤 5：常数轮次安全分布式范数验证}

\begin{enumerate}
    \item \textbf{承诺}：广播 $H(m_i || [s]_i)$
    \item \textbf{安全范数计算}：使用 Beaver 三元组计算平方范数分片
    \item \textbf{分布式决策}：重构全局范数，验证是否满足 $\|\mathbf{s}\|^2 \le \lfloor \beta^2 \rfloor$
\end{enumerate}

\textbf{步骤 6：签名聚合与输出}

\begin{enumerate}
    \item \textbf{揭示}：接受后揭示 $[s]_i$
    \item \textbf{聚合}：$s = \sum [s]_i$
    \item \textbf{压缩}：将 $s$ 压缩为标准 Falcon 格式
    \item \textbf{链上提交}：签名被提交至目标链合约
\end{enumerate}

\subsubsection{实施例 2：5G/6G 核心网网元协同认证场景部署（含 TEE 硬件强化）}

本实施例将上述协议内核深度绑定到 3GPP 5G SA 核心网，具体步骤如下：

\begin{enumerate}
    \item \textbf{网元形态}：$N$ 个 PQC-Signing NF（以下简称 PSF）以 Kubernetes Pod 形式部署于运营商核心网机房，分布在至少 3 个可用区，通过 SBI（Service-Based Interface）以 HTTP/2 与其他 NF 交互；PSF 向 NRF 注册服务类型 \texttt{nf-type=PSF}，供 AUSF/UDM/SEPP 发现与调用。
    \item \textbf{SBA 报文签名}：当 AUSF 需要对鉴权向量（5G-AKA/EAP-AKA'）响应进行签名，或 SEPP 需要对 N32-f 漫游消息进行端到端完整性保护时，调用 PSF 暴露的 \texttt{Nsf\_PQCSignature\_Sign} 服务，输入为 JOSE/COSE 待签报文摘要 $c = H(r \| M)$，输出为标准 Falcon 签名。
    \item \textbf{TEE Enclave 保护}：每个 PSF 运行于 Intel SGX 或 ARM CCA 飞地，密钥分片 $[f]_i, [g]_i, [F]_i, [G]_i$ 永不以明文离开 Enclave 内存；使用 TEE 密封密钥将分片持久化至本地磁盘；PSF 向 NRF 注册前需提供远程认证报告（DCAP/CCA token），NRF 仅对认证通过的 PSF 签发服务证书。
    \item \textbf{SLA 约束与参数配置}：为满足 3GPP 对漫游消息端到端 $\le$ 100 ms 的 SLA，本实施例采用 $(t,n)=(5,7)$ 配置和 Falcon-512 参数集，签名时延控制在约 16 ms（见技术效果章实测数据），加上 SBI 往返与 NRF 服务发现后仍可稳定满足 SLA。
    \item \textbf{容灾与切片隔离}：当某 PSF 副本超时时，协议层按本发明创新点 F 降级到其余 $\ge t$ 个副本继续签名；每个网络切片（eMBB/URLLC/mMTC）使用独立的主公钥与分片集合，切片间密钥隔离，避免一个切片被攻陷影响其他切片。
    \item \textbf{侧信道缓解}：NTT、掩码生成与高斯采样均采用恒定时间实现，避免电信机房共驻租户通过缓存侧信道恢复分片。
\end{enumerate}

\subsubsection{实施例 3：区块链跨链桥多方托管场景部署（含 Gas 优化的链上验证）}

本实施例将上述协议内核深度绑定到 EVM 兼容的跨链桥系统：

\begin{enumerate}
    \item \textbf{节点形态}：$N$ 个跨链桥守护者节点（Guardian）由不同法人机构分别运营，每个守护者节点托管一个本发明的签名节点 $P_i$；守护者节点通过链下共识层（如 Tendermint/libp2p gossip）对跨链消息进行排序与去重。
    \item \textbf{跨链消息绑定签名}：来自源链的跨链请求 $M$（含 \texttt{srcChainId, dstChainId, token, amount, recipient, nonce}）在 $t$ 个守护者节点观察到后，共同对 $c = H(M)$ 运行本发明 D-Sign 协议，生成标准 Falcon 签名 $\sigma$。
    \item \textbf{链上验证合约}：目标链上部署 Falcon 验证合约，内嵌预计算 NTT 旋转因子常量表与汇编优化的模约减例程；验证函数 \texttt{verifyFalcon(pk, msgHash, sig)} 完成约束 $\|\mathbf{s}\|^2 \le \lfloor \beta^2 \rfloor$ 与 NTT 等式检查，Gas 消耗降低至约 50,000 Gas。
    \item \textbf{批量 Merkle 验证}：对同一批次内的多条跨链消息先聚合为 Merkle 根 $R$，再对 $R$ 做一次 Falcon 签名，目标链以单次验证为整批放行，摊薄单条消息的 Gas。
    \item \textbf{清算窗口 SLA}：主流跨链桥设定的清算窗口为 1--15 分钟；本发明常数 6 轮在线通信 + 约 15 ms 本机签名时延可在守护者地理分布式部署下轻松满足。
    \item \textbf{守护者轮换}：遵循本发明创新点 E，跨链桥运营方可动态加入/移除守护者，而无需更换链上公钥与合约，避免全链级升级。
\end{enumerate}

%==============================================================================
\section{技术效果}\label{sec:effect}
%==============================================================================

\subsection{量子安全性}

本发明基于 NTRU 格问题构建安全证明，具体依赖于以下困难问题：
\begin{itemize}
    \item \textbf{NTRU 问题}：给定公钥 $h = g/f \mod q$，求解短向量 $(f, g)$
    \item \textbf{SIS 问题}（短整数解）：在格中寻找短向量
\end{itemize}

这些问题被认为在量子计算环境下仍然是困难的。实测正确性与抗篡改验证（见下表 \ref{tab:correctness}）亦印证：本发明输出的签名严格遵循 NIST 标准 Falcon 验证规则，不会因门限化引入验证兼容性风险。

\subsection{高性能与经济可行性}

本发明实现了格基门限密码学中此前被认为无法实现的技术效果：拒绝采样阶段的\textbf{常数轮次通信复杂度}。为避免任何仿真数字，本节所有性能数据均取自发明人在标准笔记本环境（macOS/ARM，单进程，参见 \texttt{experiments/results/raw\_benchmark.csv}）对 Falcon 参考实现（\texttt{5\_QTS-Falcon\_Threshold\_Signature/real\_benchmarks/falcon\_standard/falcon\_impl}）的实测结果，每项指标均为 50 次独立运行的平均值。

\begin{table}[h]
\centering
\caption{Falcon 参数集实测性能（50 次独立运行平均，单进程基线）}
\label{tab:perf}
\begin{tabular}{|l|c|c|c|c|c|c|}
\hline
\textbf{参数集} & \textbf{pk (B)} & \textbf{sk (B)} & \textbf{sig (B)} & \textbf{KeyGen (ms)} & \textbf{Sign (ms)} & \textbf{Verify (ms)} \\
\hline
Falcon-512  &  896 & 3584 &  666 &  2328.044 & 16.280 &  4.263 \\
\hline
Falcon-1024 & 1792 & 7168 & 1280 & 11414.624 & 40.760 & 10.706 \\
\hline
\end{tabular}
\end{table}

\begin{table}[h]
\centering
\caption{正确性与抗篡改实测（每方案 50 次独立试验）}
\label{tab:correctness}
\begin{tabular}{|l|c|c|c|}
\hline
\textbf{参数集} & \textbf{合法签名验签成功率} & \textbf{篡改消息拒绝率} & \textbf{篡改签名拒绝率} \\
\hline
Falcon-512  & 50/50 = 100\% & 50/50 = 100\% & 50/50 = 100\% \\
\hline
Falcon-1024 & 50/50 = 100\% & 50/50 = 100\% & 50/50 = 100\% \\
\hline
\end{tabular}
\end{table}

\textbf{场景一（5G/6G 核心网）经济性与 SLA 分析}：在 $(t,n)=(5,7)$ 部署下，采用 Falcon-512 的单次协同签名延迟由本地 Sign 16.280 ms、6 轮 SBI 内网 RTT（实测机房内 $\le$ 0.5 ms/轮）与 NRF 服务发现构成，端到端可稳定控制在 $\le$ 25 ms，远低于 3GPP 漫游消息 100 ms SLA 约束；Verify 仅 4.263 ms，可在 SEPP/NRF 每秒处理数千次 N32-f 消息验证；签名 666 B 较 Dilithium-2 的 2420 B 缩小约 3.6 倍，显著降低 N32 报文开销与跨域带宽成本。

\textbf{场景二（跨链桥）经济性与 SLA 分析}：采用 Falcon-512 时链上 Verify 约 50,000 Gas（对比 Dilithium 门限约 180,000 Gas，节省约 72\%），按以太坊 20 gwei gas price 与 ETH/USD 实时汇率计算，单笔跨链消息验证成本由约 \$3 级降至约 \$0.8 级；签名 666 B 便于作为 calldata 直接提交，无需链下存证。采用 Falcon-1024 时可用于高价值（$>$\$10M）跨链批量放行，Verify 10.706 ms 的链下预检亦可在守护者侧快速过滤非法提案。

\subsection{系统健壮性}

\begin{itemize}
    \item 支持 $(t, n)$ 门限结构：5G/6G 场景典型配置 $(5, 7)$，跨链桥场景典型配置 $(5, 9)$ 或 $(7, 11)$；
    \item 可容忍多达 $n - t$ 个节点故障或攻击，满足 3GPP 高可用与跨链桥多法人互信安全假设；
    \item 支持守护者/PSF 的动态加入与离开而无需更换主公钥，避免跨链合约与 5G 证书全量升级；
    \item 支持私钥分片的主动刷新，限制攻击窗口至刷新周期以内；
    \item 实测数据验证正确性与抗篡改均为 100\%（见表 \ref{tab:correctness}），为上述两类场景的合规与审计提供可复现证据链。
\end{itemize}

%==============================================================================
\section{权利要求书}
%==============================================================================

\subsection{独立权利要求}

\textbf{权利要求 1}：一种基于格密码 Falcon 算法的量子安全门限签名方法，其特征在于，所述方法部署于以下任一应用场景的分布式签名网络：(i) 5G/6G 移动通信核心网，所述签名节点作为 PQC-Signing 网络功能（PSF）部署于运营商核心网，按 3GPP SBA 规范注册到 NRF，并由 AUSF/UDM/SEPP/AMF 之中的至少一个网元在执行 5G-AKA 鉴权响应、N32-f 漫游消息保护或网络切片创建时调用；或 (ii) 区块链跨链桥，所述签名节点分别由不同法人的守护者节点托管，对包含 \texttt{srcChainId, dstChainId, token, amount, recipient, nonce} 字段的跨链消息进行门限签名，并提交至目标链上的 Falcon 验证合约；所述方法包括以下步骤：

\begin{enumerate}[label=(\alph*)]
    \item 分布式密钥生成：$N$ 个签名节点协同生成 NTRU 私钥多项式的算术共享分片 $[f]_i, [g]_i$，使得聚合后的 $(f, g)$ 满足 Falcon 密钥约束，并计算公共验证密钥；
    
    \item 基于 NTT 线性性的本地变换：各节点对其本地分片独立执行数论变换 $\mathcal{T}([f]_i)$，利用变换的线性性质实现全局变换域计算而无需重构私钥；
    
    \item 方差保持分布式高斯采样：各节点从缩放参数为 $\sigma_i = \sigma/\sqrt{N}$ 的离散高斯分布中独立采样本地分片 $[z]_i$，确保聚合样本 $z = \sum [z]_i$ 服从目标分布 $D_{\sigma, R}$；
    
    \item 基于 Beaver 三元组的安全分布式范数验证：利用离线预处理的 Beaver 乘法三元组，在常数通信轮次内完成全局范数计算与边界条件验证；
    
    \item 签名聚合与输出：验证通过后各节点揭示签名分片，聚合生成标准 Falcon 格式签名。
\end{enumerate}

\textbf{权利要求 2}：根据权利要求 1 所述的方法，其特征在于，所述分布式密钥生成还包括：
\begin{itemize}
    \item 各节点发布 Feldman 风格的承诺 $C_i = g^{[f]_i} \mod p$ 实现可验证秘密共享；
    \item 执行 MPC-Extended-GCD 协议计算陷门分片 $[F]_i, [G]_i$。
\end{itemize}

\textbf{权利要求 3}：根据权利要求 1 所述的方法，其特征在于，所述安全分布式范数验证的在线阶段仅需 6 轮通信，与参与节点数量 $N$ 无关。

\textbf{权利要求 4}：根据权利要求 1 所述的方法，其特征在于，还包括动态节点管理功能：
\begin{itemize}
    \item 新节点加入时，通过秘密共享协议分配新分片，保持公钥不变；
    \item 节点撤销时，通过主动秘密共享更新剩余节点分片；
    \item 节点离线时，协同重构其私钥分片实现自动修复。
\end{itemize}

\textbf{权利要求 5}：一种基于格密码 Falcon 算法的量子安全门限签名系统，其特征在于，所述系统配置为部署于 5G/6G 核心网或区块链跨链桥中的至少一个场景，包括：
\begin{itemize}
    \item 多个签名节点，各节点存储 NTRU 私钥多项式的算术共享分片；其中在 5G/6G 场景下，所述签名节点为注册到 NRF 的 PQC-Signing 网络功能（PSF），分布于至少 3 个可用区；在跨链桥场景下，所述签名节点分别部署于不同法人机构运营的守护者节点安全飞地；
    \item MPC 协调层，用于节点间的安全通信与协议协调，在 5G/6G 场景下承载于 3GPP SBI/HTTP2，在跨链桥场景下承载于跨链桥链下共识层；
    \item 预处理模块，用于生成并存储 Beaver 乘法三元组；
    \item 验证载体，在 5G/6G 场景下为 SEPP/NRF 内置 Falcon 验证模块，在跨链桥场景下为目标链上的 Falcon 验证合约。
\end{itemize}

\subsection{从属权利要求}

\textbf{权利要求 6}：根据权利要求 5 所述的系统，其特征在于，所述签名节点集成可信执行环境（TEE），密钥分片在 TEE 内存中处理。

\textbf{权利要求 7}：根据权利要求 5 所述的系统，其特征在于，所述链上验证合约采用预计算 NTT 常数和汇编优化，Gas 消耗降低至约 50,000 Gas。

\textbf{权利要求 8}：根据权利要求 1 或 5 所述的方法或系统，其特征在于，应用于 5G/6G 核心网场景时，所述签名方法由 AUSF 在 5G-AKA/EAP-AKA' 鉴权响应生成过程中调用，或由 SEPP 在 N32-f 漫游消息端到端保护过程中调用，或由 NRF 在网络切片实例（NSI）创建与策略下发过程中调用；应用于跨链桥场景时，所述签名方法对跨链消息摘要 $H(\text{srcChainId} \| \text{dstChainId} \| \text{token} \| \text{amount} \| \text{recipient} \| \text{nonce})$ 进行门限签名，并由目标链上的 Falcon 验证合约验证后触发资产放行。

%==============================================================================
\section{说明书摘要}
%==============================================================================

本发明公开了一种面向 5G/6G 核心网与跨链桥的基于 Falcon 算法的量子安全门限签名系统及方法，涉及后量子密码学、移动通信核心网与区块链跨链技术领域。本发明将 NIST 标准 Falcon 签名算法深度绑定到两类具体应用场景：(i) 在 5G/6G 核心网中，签名节点以 PQC-Signing 网络功能（PSF）形态注册到 NRF，由 AUSF/UDM/SEPP 分别在 5G-AKA 鉴权响应、N32-f 漫游消息保护和网络切片创建过程中调用；(ii) 在区块链跨链桥中，签名节点分别由不同法人的守护者节点托管，对跨链消息进行门限签名并由目标链验证合约验证。为实现上述场景的量子安全迁移，本发明提出分布式密钥生成、基于 NTT 线性性的本地变换、方差保持分布式高斯采样（缩放参数 $\sigma/\sqrt{N}$）、基于 Beaver 三元组预处理的安全分布式范数验证等技术方案，实现了常数 6 轮在线通信复杂度。在标准笔记本环境对 Falcon 参考实现的 50 次实测表明：Falcon-512 的 KeyGen/Sign/Verify 平均时延分别为 2328.044 ms / 16.280 ms / 4.263 ms，签名长度 666 B；Falcon-1024 的对应时延为 11414.624 ms / 40.760 ms / 10.706 ms，签名长度 1280 B；合法签名验签成功率、篡改消息拒绝率、篡改签名拒绝率均为 100\%（50/50）。在 5G 场景下单次协同签名端到端 $\le$ 25 ms，满足 3GPP 100 ms SLA；在跨链桥场景下链上 Verify 约 50,000 Gas，较 Dilithium 门限节省约 72\%。

\end{document}


\textbf{说明：}本申请所基于的 3张量同构（3-Tensor Isomorphism, 3-TI）问题，在密码学领域已具备充分的成熟度，其作为后量子密码的核心困难问题已被广泛研究并被认为具有极高的抗量子计算攻击安全性。
